%%
%% RGSW, External Product and Internal Product
%%
\section{TFHE Bootstrapping Primitives}
\label{sec:tfhe}
% Scope: the RGSW machinery this thesis builds on --- RGSW ciphertexts, the
% external and internal products, and the generalization to arbitrary gadgets.
% Bootstrapping itself is only referenced at a high level (see FHE section).

TFHE~\cite{cryptoeprint:2016/870, cryptoeprint:2018/421} is optimized for computations of large multiplicative depth on small (binary) plaintexts, and is best known for its fast bootstrapping procedure. This section covers only the machinery that bootstrapping --- and this thesis --- is built from: RGSW ciphertexts and the external and internal products. Bootstrapping itself (blind rotation, sample extraction, LWE key switching) is beyond the scope of this work.

GSW encryption was introduced over plain LWE by Gentry, Sahai and Waters~\cite{cryptoeprint:2013/340}. The ring variant (RGSW) appears in FHEW~\cite{cryptoeprint:2014/816}, which uses RGSW accumulators for bootstrapping, and in SHIELD~\cite{cryptoeprint:2014/838} around the same time. The external product --- the asymmetric RGSW $\times$ RLWE multiplication, as opposed to GSW's original symmetric product --- is due to Chillotti, Gama, Georgieva and Izabach\`ene~\cite{cryptoeprint:2016/870} (journal version~\cite{cryptoeprint:2018/421}), with a concurrent equivalent (the ``hybrid product'') by Brakerski and Perlman~\cite{cryptoeprint:2016/339}. The internal product (RGSW $\times$ RGSW) is the ring version of GSW's original multiplication.

\begin{remark}[The torus pivot, once and for all]
    TFHE is natively formulated over the (discretized) torus, with TLWE/TGSW in place of RLWE/RGSW. The discretized torus $q^{-1}\mathbb{Z}/\mathbb{Z}$ is isomorphic to $\mathbb{Z}_q$ via $x \mapsto qx$~\cite{cryptoeprint:2021/1402}, under which the torus gadget entries $1/B^j$ become $q/B^j$. All statements transfer verbatim; we use ring notation over $R_Q$ throughout, with RLWE ciphertexts as in \autoref{sec:fhe} and gadget pairs $(D, P)$ satisfying the (approximate) gadget property of Def.~\ref{def:gadget}.
\end{remark}

\subsection{RGSW Ciphertexts}
Fix the radix gadget of Def.~\ref{ex:radix_gadget}, $g = P(1) = (q/B, \dots, q/B^\ell)^T$, and the corresponding \textit{gadget matrix}
\begin{equation}
    G = I_2 \otimes g \in R_Q^{2\ell \times 2}.
\end{equation}

\begin{definition}[RGSW Ciphertext]
    An RGSW encryption of $\mu \in R$ is
    \begin{equation}
        \label{def:rgsw}
        C = Z + \mu\cdot G \in R_Q^{2\ell \times 2},
        \qquad \text{each row of } Z \text{ an RLWE encryption of } 0.
    \end{equation}
\end{definition}

Structurally, an RGSW ciphertext is $2\ell$ RLWE ciphertexts encrypting gadget-scaled copies of $\mu$: the rows of one block encrypt $\mu\cdot g_j$ directly, and the rows of the other encrypt $-s\cdot\mu g_j$. In particular, RGSW is additively homomorphic (matrix addition), and $\mu$ can be recovered by decrypting any row of the $\mu\cdot g_j$ block.

\subsection{External Product}
Gadget decomposition extends to an RLWE ciphertext $c = (c_0, c_1)$ by decomposing each component:
\begin{equation}
    G^{-1}(c) = \big(D(c_0) \,\big|\, D(c_1)\big) \in R^{2\ell},
    \qquad \text{satisfying } G^{-1}(c)\cdot G = (c_0, c_1) + (e_0, e_1),
    \quad \|e_i\|_\infty \le \varepsilon.
\end{equation}
% exact gadgets: ε = 0; radix gadget over torus: ε = 1/(2B^ℓ) scaled

\begin{definition}[External Product]
    The external product $\boxdot$ of an RGSW ciphertext $C$ and an RLWE ciphertext $c$ is~\cite{cryptoeprint:2016/870, cryptoeprint:2016/339}
    \begin{equation}
        \label{eq:ext_prod}
        C \boxdot c = G^{-1}(c) \cdot C \in R_Q^2.
    \end{equation}
\end{definition}

Correctness follows in one line by expanding \eqref{def:rgsw}:
\begin{equation*}
    G^{-1}(c)\cdot(Z + \mu_C G)
    = \underbrace{G^{-1}(c)\cdot Z}_{\text{decomposition--noise term}}
    + \;\mu_C\cdot\big(c + e_{\mathrm{dec}}\big)
    \;=\; \rlwe(\mu_C \cdot \mu_c) + \text{noise},
\end{equation*}
i.e.\ the result encrypts the product of the two messages, with worst-case noise
\begin{equation}
    \label{eq:ext_prod_noise}
    \|v_{\boxdot}\|_\infty \;\lesssim\;
    \underbrace{2\ell\, N\, \beta\, B_{Z}}_{\langle D(c_i),\, \vec{e}_Z\rangle}
    \;+\;
    \underbrace{\|\mu_C\|_1 \cdot \varepsilon}_{\mu_C \cdot e_{\mathrm{dec}}}
    \;+\; \|\mu_C\|_1 \cdot \|v_c\|_\infty,
\end{equation}
where $\beta$ is the digit bound of $D$, $\varepsilon$ its precision, $B_Z$ the RGSW noise and $v_c$ the noise of the RLWE input (statement only; proof deferred to Thm.~\ref{thm:generalization}).

Two consequences are worth flagging. First, the noise is \emph{asymmetric}: the RLWE input's noise enters scaled by $\mu_C$, while the RGSW noise enters only through small digits. The RGSW message must therefore be small --- typically $\mu_C \in \{0,1\}$ or a monomial $X^k$ --- while $\mu_c$ is unconstrained. With $\|\mu_C\|_1 = 1$ the input noise passes through \emph{unchanged}, so a chain of $n$ external products grows the noise only linearly in $n$; this near-constant growth per multiplication is what TFHE bootstrapping --- and the \gls{spar} protocol --- are built on.
% this asymmetry is why CMux / tree traversal chains external products on the RLWE side
Second, the cost is $2\ell$ ring multiplications plus the cost of the decomposition $D$.

\subsection{Internal Product}
The internal product $\boxtimes$ (RGSW $\times$ RGSW $\to$ RGSW) applies the external product row-wise, each row of $C_1$ being an RLWE ciphertext:
\begin{equation}
    C_1 \boxtimes C_2 = \big(C_2 \boxdot \mathrm{row}_r(C_1)\big)_{r=1}^{2\ell}.
\end{equation}
It costs $2\ell$ external products, with the noise of \eqref{eq:ext_prod_noise} per row. Its role is composing RGSW ciphertexts (e.g.\ products of control bits); it is not used on the critical path of this thesis and is included for completeness.

\subsection{Generalization}
Nothing above used the radix structure of $g$ --- only the gadget property of $(D, P)$. Any pair from \autoref{sec:gadgets} therefore yields a valid RGSW ciphertext and external product.

\begin{theorem}
    \label{thm:generalization}
    Let $D(a), P(b) \in R_Q^\ell$ satisfy the approximate gadget property
    (Def.~\ref{def:gadget}) with digit bound $\beta$ and decomposition error
    $\varepsilon$. Then
    \begin{align}
        G &= I_2 \otimes P(1) \label{eq:gen_G}\\
        G^{-1}(x) &= \big(D(x_0) \,\big|\, D(x_1)\big) \label{eq:gen_Ginv}
    \end{align}
    define a valid RGSW ciphertext \eqref{def:rgsw} and external product
    \eqref{eq:ext_prod} with
    $\mathrm{RGSW}(\mu_C) \boxdot \mathrm{RLWE}(\mu_c)
     = \mathrm{RLWE}(\mu_C \cdot \mu_c)$,
    provided the resulting noise remains below the decryption bound.
    % Proof sketch: (i) the inexact gadget identity G^{-1}(x)·G = x + e holds
    % componentwise from Def.~\ref{def:gadget}; (ii) rerun the external-product
    % expansion; noise bound with β and ε as displayed above.
\end{theorem}

Instantiating Thm.~\ref{thm:generalization} with the gadgets of \autoref{sec:gadgets} gives external products computable natively in RNS: the BV-type gadgets are exact ($\varepsilon = 0$), while the extension-based gadget (GHS/Hybrid, \autoref{sec:gadgets:ghs}) is approximate, with $\varepsilon > 0$ from the FastBConv overflow and the ApproxModDown rounding, and the $P$-scaling absorbed into $P(\cdot)$.
% this bridge sentence is the load-bearing edge of the related-works graph

Finally, a genuine (not merely notational) difference from the TFHE setting: in BGV the zero-encryptions carry the plaintext factor $t$ --- $Z$-rows of the form $(as + te, a)$ --- so the noise bounds above acquire factors of $t$. This is treated in \autoref{chp:methodology}.

\newpage
